

%%%%%%%%%%%%%%%%%%%%%%%%%%%%%%%%%
\subsection{Ray tracing simulations}
%%%%%%%%%%%%%%%%%%%%%%%%%%%%%%%%%%

\textbf{You will:}

\begin{itemize}
    \item learn to define geometrical sources
    \item understand the use of a python script for modifying an existing source
    \item understand reference frames
\end{itemize}

\textbf{Questions:}

\textbf{i)} Create a collimated (i.e., zero divergence) geometric source with rectangular shape with vertical length 2 mm (X coordinate, column 1) and horizontal length 1 mm (Z coordinate, column 3).

\textbf{ii)} Shift the beam (1, 0, 0.5) mm to get all rays positive.

\textbf{iii)} Create an empty optical element, with incident angle 0 deg (normal incidence) and reflection angle 180$^\circ$ (the same direction), and p=q=0. Trace the system fr optical element orientation angles $\alpha$ equal to 0, 90, 180 and 270 degrees. Verify the results with the pictures shown before.

\textbf{Hints:} you may load the workspace {\tt reference\_frame.ows}. 

\textbf{Answer}

\begin{figure}[H]
\label{fig:fig1}
\caption{Effect of an empty element in the beam for different $\alpha$. Note that the beam is unchanged, but the axes are rotated and displayed in the same configuration, therefore there is an apparent rotation of the beam.}
\includegraphics[width=0.9\textwidth]{reference-frames/fig1.png}
\end{figure}
